% !TeX root = main.tex

% Два размера рисунков, как в лабораторной работе № 3.
% Аргументы: имя файла в images, подпись, метка для ссылки.
% Если файл ещё не добавлен, вместо него выводится рамка высотой 3 см.
\newcommand{\imageplaceholder}[3]{%
  \begin{figure}[H]
    \centering
    \IfFileExists{images/#1}{%
      \includegraphics[width=0.95\textwidth]{#1}%
    }{%
      \fbox{%
        \parbox[c][3cm][c]{0.95\textwidth}{%
          \centering
          Место для рисунка\\[0.4em]
          \texttt{images/\detokenize{#1}}
        }%
      }%
    }
    \caption{#2}
    \label{#3}
  \end{figure}
}


\tableofcontents
\newpage

\section{Цель работы}

Освоить принципы построения робастной системы управления линейным
многомерным объектом по состоянию на основе метода функций Ляпунова.

\section{Постановка задачи и исходные данные}

Рассматривается возмущённый объект управления
\begin{equation}
  \dot{\bm x}=A\bm x+\bm b u+\bm\delta(t),
  \qquad y=C\bm x,
  \qquad \bm x\in\mathbb R^2.
  \label{eq:plant}
\end{equation}
В данной работе используются матрицы объекта и сигнал задания из работы
№ 3. Для варианта 1 согласно таблице 3.1 методических указаний (с.~34)
заданы
\begin{equation}
  A=\begin{pmatrix}0&1\\1&1\end{pmatrix},
  \qquad
  \bm b=\begin{pmatrix}0\\1\end{pmatrix},
  \qquad C=\begin{pmatrix}1&0\end{pmatrix},
  \qquad b_0=1.
\end{equation}
Задающее воздействие имеет вид
\begin{equation}
  g(t)=3\sgn\bigl(\cos(0{,}2t)\bigr)+3.
\end{equation}

В опытах с возмущением
используется вектор-функция
\begin{equation}
  \bm\delta(t)=
  \begin{pmatrix}
    0{,}6\sin(10t)+0{,}1\sin(50t)\\
    0{,}5\cos(12t)+0{,}2\sin(30t)
  \end{pmatrix}.
  \label{eq:disturbance}
\end{equation}
В опытах без возмущения принимается $\bm\delta(t)\equiv\bm 0$.

Ошибка слежения и цель управления определяются соотношениями
\begin{equation}
  \bm e(t)=\bm x_M(t)-\bm x(t),
  \qquad \norm{\bm e(t)}\le\Delta
  \quad\text{при всех }t\ge T,
  \label{eq:control-goal}
\end{equation}
где $\Delta>0$ \textemdash{} допустимая установившаяся ошибка,
$T$ \textemdash{} время настройки, $\bm x_M$ \textemdash{} состояние
эталонной модели.

\section{Построение моделей систем управления}

\subsection{Эталонная модель и общие параметры}

Эталонная модель сохраняет структуру, использованную в лабораторной
работе № 3:
\begin{equation}
  \dot{\bm x}_M=A_M\bm x_M+\bm b_M g,
  \qquad y_M=C\bm x_M,
\end{equation}
\begin{equation}
  A_M=\begin{pmatrix}0&1\\-a_{M0}&-a_{M1}\end{pmatrix},
  \qquad \bm b_M=\begin{pmatrix}0\\a_{M0}\end{pmatrix}.
\end{equation}
Условия согласования и общий закон управления:
\begin{equation}
  A_M=A+\bm b\bm\theta^{\mathsf T},
  \qquad \bm b=\kappa\bm b_M,
  \qquad u=\hat{\bm\theta}^{\mathsf T}\bm x+\frac{1}{\kappa}g.
  \label{eq:control}
\end{equation}
Матрица $P=P^{\mathsf T}>0$ определяется из уравнения Сильвестра
\begin{equation}
  A_M^{\mathsf T}P+PA_M=-Q,
  \qquad Q=Q^{\mathsf T}>0.
\end{equation}
Для построения графиков параметрических ошибок используется
\begin{equation}
  \tilde{\bm\theta}=\bm\theta-\hat{\bm\theta}.
\end{equation}

Матрицы синтезированной в работе №3 эталонной модели имеют вид
\begin{equation}
  A_M=\begin{pmatrix}
    0&1\\
    -1329{,}4782&-72{,}9240
  \end{pmatrix},
  \qquad
  \bm b_M=\begin{pmatrix}
    0\\
    1329{,}4782
  \end{pmatrix},
  \qquad C_M=\begin{pmatrix}1&0\end{pmatrix}
\end{equation}

\begin{equation}
  Q=I,
  \qquad \kappa=\frac{1}{1329{,}4782}\approx0{,}0007521748,
  \qquad P\approx\begin{pmatrix}
    9{,}1497874344&0{,}0003760874\\
    0{,}0003760874&0{,}0068616105
  \end{pmatrix}
\end{equation}

\begin{equation}
  \bm\theta=
  \begin{pmatrix}
    -1330{,}4782\\
    -73{,}9240
  \end{pmatrix}
\end{equation}

\subsection{Нелинейный робастный регулятор}

В первом решении используется закон управления \eqref{eq:control}
со статической настройкой параметров по формуле (4.4) методических указаний:
\begin{equation}
  \hat{\bm\theta}=\gamma\bm x\bm b^{\mathsf T}P\bm e,
  \qquad \gamma>0.
  \label{eq:static-law}
\end{equation}
Коэффициент $\gamma$ определяет интенсивность нелинейной обратной связи.

\subsection{Адаптивный регулятор с параметрической обратной связью}

Во втором решении закон управления \eqref{eq:control} дополняется
алгоритмом (4.6) методических указаний:
\begin{equation}
  \dot{\hat{\bm\theta}}
  =-\sigma\hat{\bm\theta}+\gamma\bm x\bm b^{\mathsf T}P\bm e,
  \qquad \gamma>0,
  \qquad \sigma>0.
  \label{eq:adaptive-law}
\end{equation}
Здесь $\gamma$ \textemdash{} коэффициент адаптации,
$\sigma$ \textemdash{} коэффициент параметрической обратной связи.


\section{Исследование нелинейного робастного управления}

Исследуется система
\eqref{eq:control}, \eqref{eq:static-law} при трёх различных значениях
$\gamma$. Для каждого значения выполняются два опыта: без возмущения
и с возмущением \eqref{eq:disturbance}. В каждом опыте траектории
$\bm x(t)$ и $\bm x_M(t)$ приводятся на одном графике, а компоненты
ошибки $\bm e(t)$ \textemdash{} на другом.

Выберем $\gamma=\{10000, 100000, 1000000\}$.

\subsection{Первое значение коэффициента обратной связи}

% TODO: указать численное значение gamma_1.
При $\gamma=10000$ результаты опыта без возмущения предусмотрены
на рисунках~\ref{fig:robust-g2-free-x}, \ref{fig:robust-g2-free-e}.

\imageplaceholder{robust_g2_free_x.png}
  {Траектории объекта и модели при $10000$, без возмущения}
  {fig:robust-g2-free-x}
\imageplaceholder{robust_g2_free_e.png}
  {Ошибка слежения при $10000$, без возмущения}
  {fig:robust-g2-free-e}

Результаты опыта с возмущением предусмотрены
на рисунках~\ref{fig:robust-g2-dist-x}, \ref{fig:robust-g2-dist-e}.

\imageplaceholder{robust_g2_dist_x.png}
  {Траектории объекта и модели при $10000$, с возмущением}
  {fig:robust-g2-dist-x}
\imageplaceholder{robust_g2_dist_e.png}
  {Ошибка слежения при $10000$, с возмущением}
  {fig:robust-g2-dist-e}

\subsection{Второе значение коэффициента обратной связи}

% TODO: указать численное значение gamma_2.
При $\gamma=100000$ результаты опыта без возмущения предусмотрены
на рисунках~\ref{fig:robust-g3-free-x}, \ref{fig:robust-g3-free-e}.

\imageplaceholder{robust_g3_free_x.png}
  {Траектории объекта и модели при $100000$, без возмущения}
  {fig:robust-g3-free-x}
\imageplaceholder{robust_g3_free_e.png}
  {Ошибка слежения при $100000$, без возмущения}
  {fig:robust-g3-free-e}

Результаты опыта с возмущением предусмотрены
на рисунках~\ref{fig:robust-g3-dist-x}, \ref{fig:robust-g3-dist-e}.

\imageplaceholder{robust_g3_dist_x.png}
  {Траектории объекта и модели при $100000$, с возмущением}
  {fig:robust-g3-dist-x}
\imageplaceholder{robust_g3_dist_e.png}
  {Ошибка слежения при $100000$, с возмущением}
  {fig:robust-g3-dist-e}

\subsection{Третье значение коэффициента обратной связи}

% TODO: указать численное значение gamma_3.
При $\gamma=1000000$ результаты опыта без возмущения предусмотрены
на рисунках~\ref{fig:robust-g3-free-x}, \ref{fig:robust-g3-free-e}.

\imageplaceholder{robust_g3_free_x.png}
  {Траектории объекта и модели при $1000000$, без возмущения}
  {fig:robust-g3-free-x}
\imageplaceholder{robust_g3_free_e.png}
  {Ошибка слежения при $1000000$, без возмущения}
  {fig:robust-g3-free-e}

Результаты опыта с возмущением предусмотрены
на рисунках~\ref{fig:robust-g3-dist-x}, \ref{fig:robust-g3-dist-e}.

\imageplaceholder{robust_g3_dist_x.png}
  {Траектории объекта и модели при $1000000$, с возмущением}
  {fig:robust-g3-dist-x}
\imageplaceholder{robust_g3_dist_e.png}
  {Ошибка слежения при $1000000$, с возмущением}
  {fig:robust-g3-dist-e}

\subsection{Анализ результатов}

По результатам моделирования установлено, что увеличение коэффициента нелинейной обратной связи \(\gamma\) приводит к существенному уменьшению ошибки слежения \(e(t)=x_M(t)-x(t)\). При \(\gamma=10^4\) наблюдаются значительные отклонения траекторий объекта от эталонной модели, тогда как при \(\gamma=10^6\) обеспечивается значительно более точное воспроизведение заданной траектории.

При наличии внешнего возмущения ошибка приобретает дополнительные колебания, однако остаётся ограниченной. С увеличением \(\gamma\) влияние возмущения на качество слежения уменьшается. При этом даже в отсутствие возмущения ошибка не стремится к нулю, что соответствует теоретическим свойствам статического робастного регулятора (4.4).

\noindent\textbf{Вывод по пункту:} Таким образом, результаты моделирования подтверждают работоспособность нелинейного робастного управления и возможность повышения точности слежения путём увеличения коэффициента \(\gamma\).


\section{Исследование адаптивного и робастного управления}

По пункту 3 порядка выполнения работы исследуется система
\eqref{eq:control}, \eqref{eq:adaptive-law}. При $\sigma=\sigma_1$
сравниваются два значения $\gamma$ из предыдущего раздела.
Затем для одного из них, обозначенного $\gamma_*$, опыт повторяется
при уменьшенном коэффициенте $\sigma_2<\sigma_1$.
Каждая комбинация проверяется без возмущения и с возмущением.
Для каждого опыта предусмотрены три графика: $\bm x(t)$ и $\bm x_M(t)$
на одном поле, ошибка слежения $\bm e(t)$ и параметрическая ошибка
$\tilde{\bm\theta}(t)$ на отдельных полях.

\subsection{Первое значение коэффициента адаптации}

% TODO: указать численные значения sigma_1 и gamma_1.
При $\sigma=\sigma_1$, $\gamma=100000$ результаты опыта без возмущения
предусмотрены на рисунках~\mbox{\ref{fig:adaptive-s1-g2-free-x}\textendash\ref{fig:adaptive-s1-g2-free-theta}}.

\imageplaceholder{adaptive_s1_g2_free_x.png}
  {Траектории при $\sigma_1$, $100000$, без возмущения}
  {fig:adaptive-s1-g2-free-x}
\imageplaceholder{adaptive_s1_g2_free_e.png}
  {Ошибка слежения при $\sigma_1$, $100000$, без возмущения}
  {fig:adaptive-s1-g2-free-e}
\imageplaceholder{adaptive_s1_g2_free_theta.png}
  {Параметрическая ошибка при $\sigma_1$, $100000$, без возмущения}
  {fig:adaptive-s1-g2-free-theta}

Результаты опыта с возмущением предусмотрены
на рисунках~\mbox{\ref{fig:adaptive-s1-g2-dist-x}\textendash\ref{fig:adaptive-s1-g2-dist-theta}}.

\imageplaceholder{adaptive_s1_g2_dist_x.png}
  {Траектории при $\sigma_1$, $100000$, с возмущением}
  {fig:adaptive-s1-g2-dist-x}
\imageplaceholder{adaptive_s1_g2_dist_e.png}
  {Ошибка слежения при $\sigma_1$, $100000$, с возмущением}
  {fig:adaptive-s1-g2-dist-e}
\imageplaceholder{adaptive_s1_g2_dist_theta.png}
  {Параметрическая ошибка при $\sigma_1$, $100000$, с возмущением}
  {fig:adaptive-s1-g2-dist-theta}

\subsection{Второе значение коэффициента адаптации}

% TODO: указать численное значение gamma_2; sigma_1 сохранить.
При $\sigma=\sigma_1$, $\gamma=1000000$ результаты опыта без возмущения
предусмотрены на рисунках~\mbox{\ref{fig:adaptive-s1-g3-free-x}\textendash\ref{fig:adaptive-s1-g3-free-theta}}.

\imageplaceholder{adaptive_s1_g3_free_x.png}
  {Траектории при $\sigma_1$, $1000000$, без возмущения}
  {fig:adaptive-s1-g3-free-x}
\imageplaceholder{adaptive_s1_g3_free_e.png}
  {Ошибка слежения при $\sigma_1$, $1000000$, без возмущения}
  {fig:adaptive-s1-g3-free-e}
\imageplaceholder{adaptive_s1_g3_free_theta.png}
  {Параметрическая ошибка при $\sigma_1$, $1000000$, без возмущения}
  {fig:adaptive-s1-g3-free-theta}

Результаты опыта с возмущением предусмотрены
на рисунках~\mbox{\ref{fig:adaptive-s1-g3-dist-x}\textendash\ref{fig:adaptive-s1-g3-dist-theta}}.

\imageplaceholder{adaptive_s1_g3_dist_x.png}
  {Траектории при $\sigma_1$, $1000000$, с возмущением}
  {fig:adaptive-s1-g3-dist-x}
\imageplaceholder{adaptive_s1_g3_dist_e.png}
  {Ошибка слежения при $\sigma_1$, $1000000$, с возмущением}
  {fig:adaptive-s1-g3-dist-e}
\imageplaceholder{adaptive_s1_g3_dist_theta.png}
  {Параметрическая ошибка при $\sigma_1$, $1000000$, с возмущением}
  {fig:adaptive-s1-g3-dist-theta}

\subsection{Уменьшение коэффициента параметрической обратной связи}

% TODO: указать sigma_2 < sigma_1 и выбранное gamma_* из gamma_1, gamma_2.
При $\sigma=\sigma_2<\sigma_1$, $\gamma=\gamma_*$ результаты опыта
без возмущения предусмотрены
на рисунках~\mbox{\ref{fig:adaptive-s2-free-x}\textendash\ref{fig:adaptive-s2-free-theta}}.
Остальные условия сохраняются такими же, как в соответствующем опыте
при $\sigma_1$.

\imageplaceholder{adaptive_s2_free_x.png}
  {Траектории при $\sigma_2$, $\gamma_*$, без возмущения}
  {fig:adaptive-s2-free-x}
\imageplaceholder{adaptive_s2_free_e.png}
  {Ошибка слежения при $\sigma_2$, $\gamma_*$, без возмущения}
  {fig:adaptive-s2-free-e}
\imageplaceholder{adaptive_s2_free_theta.png}
  {Параметрическая ошибка при $\sigma_2$, $\gamma_*$, без возмущения}
  {fig:adaptive-s2-free-theta}

Результаты опыта с возмущением предусмотрены
на рисунках~\mbox{\ref{fig:adaptive-s2-dist-x}\textendash\ref{fig:adaptive-s2-dist-theta}}.

\imageplaceholder{adaptive_s2_dist_x.png}
  {Траектории при $\sigma_2$, $\gamma_*$, с возмущением}
  {fig:adaptive-s2-dist-x}
\imageplaceholder{adaptive_s2_dist_e.png}
  {Ошибка слежения при $\sigma_2$, $\gamma_*$, с возмущением}
  {fig:adaptive-s2-dist-e}
\imageplaceholder{adaptive_s2_dist_theta.png}
  {Параметрическая ошибка при $\sigma_2$, $\gamma_*$, с возмущением}
  {fig:adaptive-s2-dist-theta}

\subsection{Анализ результатов}

\textit{После моделирования сравнить ошибки слежения и параметрические
ошибки при изменении $\gamma$ и $\sigma$. Оценить влияние возмущения,
остаточную ошибку и ограниченность сигналов. Сопоставить результаты
со статическим робастным регулятором.}

\noindent\textbf{Вывод по пункту:} \textit{Заполнить по результатам опытов.}

\section{Выводы}

\textit{Сформулировать выводы по результатам исследования двух алгоритмов:
влияние коэффициентов $\gamma$ и $\sigma$, влияние возмущения,
достигаемая точность слежения и различия между регуляторами.}
